\paragraph{扫描谱指纹的 Carath\'eodory--Schur 提升与收缩实现：右半平面正性、最小实现与局部判别式链}
沿用结论 \ref{con:xi-2kappa-minimal-complete-samples} 的一侧谱指纹表示：
存在 \(\kappa\) 个参数 \(r_j\in\mathbb{D}\) 与权重 \(c_j>0\) 使得
\begin{equation}\label{eq:xi-scan-fingerprint-one-sided-exponential-sum}
\boxed{\;
u_n=\sum_{j=1}^{\kappa}c_j\,r_j^n\qquad(n\ge 0)\;}
\end{equation}
（在一般位置下该表示取为最小：\(\{r_j\}\) 两两不同）。

\begin{proposition}[扫描谱指纹的 Carath\'eodory 提升：右半平面正性与系数可逆性]\label{prop:xi-scan-fingerprint-caratheodory-lift}
\leanverified{paper\_xi\_scan\_fingerprint\_caratheodory\_lift}
在 \eqref{eq:xi-scan-fingerprint-one-sided-exponential-sum} 的设定下，定义解析函数
\begin{equation}\label{eq:xi-scan-fingerprint-carath}
\mathscr{C}_u(z):=\sum_{j=1}^{\kappa}c_j\,\frac{1+r_j z}{1-r_j z},
\qquad |z|<1.
\end{equation}
则对一切 \(|z|<1\) 有严格正性 \(\Re \mathscr{C}_u(z)>0\)。
并且 \(\mathscr{C}_u\) 的 Taylor 系数与 \((u_n)\) 一一对应：
\begin{equation}\label{eq:xi-scan-fingerprint-carath-series}
\boxed{\;
\mathscr{C}_u(z)
=u_0+2\sum_{n\ge 1}u_n\,z^n,\qquad u_0=\sum_{j=1}^{\kappa}c_j>0.\;}
\end{equation}
因此有限原子扫描谱指纹与“有限极点 Carath\'eodory 函数”严格等价：\((u_n)_{n\ge 0}\) 在制度可见域内被唯一封口为一条右半平面正性轨道 \(\mathscr{C}_u\)。
\end{proposition}
\begin{proof}
令 \(w=r_j z\)。由于 \(|r_j|<1\) 且 \(|z|<1\)，有 \(|w|<1\)。
而 Cayley 变换满足
\[
\Re\frac{1+w}{1-w}=\frac{1-|w|^2}{|1-w|^2}>0,
\]
故 \(\Re\mathscr{C}_u(z)=\sum_j c_j\,\Re\frac{1+r_j z}{1-r_j z}>0\)。
展开
\(
\frac{1+w}{1-w}=1+2\sum_{n\ge 1}w^n
\)
并代入 \eqref{eq:xi-scan-fingerprint-one-sided-exponential-sum} 即得 \eqref{eq:xi-scan-fingerprint-carath-series}。证毕。
\end{proof}

\begin{theorem}[Pick 核的规范刚性与 strictification 商结构]\label{thm:xi-pick-kernel-gauge-rigidity}
\leanverified{paper\_xi\_pick\_kernel\_gauge\_rigidity}
令 \(C_1,C_2\) 为单位圆盘 \(\mathbb D\) 上的 Carath\'eodory 函数（解析且 \(\Re C_j>0\)）。
定义对应的 Pick 核
\[
K_{C_j}(w,z):=\frac{C_j(w)+\overline{C_j(z)}}{1-w\overline z}\qquad(w,z\in\mathbb D).
\]
若 \(K_{C_1}\equiv K_{C_2}\) 于 \(\mathbb D\times\mathbb D\)，则存在且仅存在 \(\eta\in\RR\) 使
\[
C_1(w)\equiv C_2(w)+\mathrm{i}\eta.
\]
因此，一切仅依赖 Pick 核 \(K_C\) 的构造（例如 Toeplitz--PSD 塔、Schur/CMV 参数、Schur--Cohn LMI 等）都严格因子化为 Carath\'eodory 类空间的商
\[
\mathrm{Car}(\mathbb D)/(\mathrm{i}\RR),
\]
并且严格化（strictification）\(C\mapsto C+\mathrm{i}\eta\)（不改变核）给出该核层级上的全部规范自由度。
\end{theorem}
\begin{proof}
由 \(K_{C_1}\equiv K_{C_2}\) 得
\[
C_1(w)+\overline{C_1(z)}\equiv C_2(w)+\overline{C_2(z)}\qquad(w,z\in\mathbb D).
\]
取 \(z=0\) 得
\(
C_1(w)-C_2(w)\equiv \overline{C_2(0)-C_1(0)}=:\lambda
\)
为常数。
再取 \(w=0\) 得
\[
C_1(0)+\overline{C_1(0)}=C_2(0)+\overline{C_2(0)}
\quad\Longrightarrow\quad
\Re\lambda=0,
\]
故 \(\lambda=\mathrm{i}\eta\) 且 \(\eta\in\RR\)。
唯一性由常数差的确定性直接推出。证毕。
\end{proof}

\begin{proposition}[最小收缩实现的唯一性：谱指纹即一维缺陷收缩的状态矩]\label{prop:xi-scan-fingerprint-minimal-contraction-realization}
令 \(\mathscr{C}_u\) 如 \eqref{eq:xi-scan-fingerprint-carath}，并定义 Schur 函数
\[
s(z):=\frac{\mathscr{C}_u(z)-\mathscr{C}_u(0)}{\mathscr{C}_u(z)+\mathscr{C}_u(0)},
\qquad |z|<1.
\]
则 \(s\) 属于 Schur 类：\(|s(z)|<1\) 且 \(s(0)=0\)。
并且存在一个（复）Hilbert 空间 \(\mathcal{H}\) 上的完全非幺正收缩算子 \(T\) 与单位向量 \(e\in\mathcal{H}\) 使得对一切 \(|z|<1\) 有
\begin{equation}\label{eq:xi-scan-fingerprint-contraction-herglotz}
\boxed{\;
\mathscr{C}_u(z)=\mathscr{C}_u(0)\,\langle e,(I+zT)(I-zT)^{-1}e\rangle\;}
\end{equation}
从而得到谱矩恒等式
\begin{equation}\label{eq:xi-scan-fingerprint-contraction-moments}
\boxed{\;
\frac{u_n}{u_0}=\langle e,T^{n}e\rangle\qquad(n\ge 0).\;}
\end{equation}
在最小性条件
\(
\overline{\Span}\{T^n e:n\ge 0\}=\mathcal{H}
\)
下，该实现 \((T,e)\) 在幺正等价意义下唯一。
在有限原子情形（\(\mathscr{C}_u\) 有理且表示最小）下，可取 \(\dim\mathcal{H}=\kappa\)，并且 \(\mathrm{Spec}(T)=\{r_j\}_{j=1}^{\kappa}\)（按代数重数计）。
\end{proposition}
\begin{proof}[证明要点]
由 \(\Re\mathscr{C}_u>0\) 与 \(\mathscr{C}_u(0)=u_0>0\)，Cayley 变换给出 \(|s(z)|<1\)。
\par
对存在性，在有限原子表示下取 \(\mathcal{H}=\CC^\kappa\)，令 \(T:=\mathrm{diag}(r_1,\dots,r_\kappa)\)，并令
\(
e:=(\sqrt{c_1/u_0},\dots,\sqrt{c_\kappa/u_0})^\top
\)。
则 \(\|e\|=1\) 且
\[
u_0\,\langle e,(I+zT)(I-zT)^{-1}e\rangle
=\sum_{j=1}^{\kappa}c_j\,\frac{1+r_j z}{1-r_j z}
=\mathscr{C}_u(z),
\]
即得 \eqref{eq:xi-scan-fingerprint-contraction-herglotz}。
将 \((I+zT)(I-zT)^{-1}=I+2\sum_{n\ge 1}z^nT^n\) 代入并与 \eqref{eq:xi-scan-fingerprint-carath-series} 比较，得到 \eqref{eq:xi-scan-fingerprint-contraction-moments}。
最小表示下 \(e\) 在每个特征向量方向的投影非零，故循环子空间为全空间。
\par
唯一性属于最小收缩 realization 的标准结论：若两组最小实现给出同一 \(\mathscr{C}_u\)（等价地同一 \((u_n)\)），则其由矩序列 \(\langle e,T^n e\rangle\) 所确定的 cyclic 表示在幺正意义下同构。证毕。
\end{proof}

\begin{corollary}[McMillan 度判别：\texorpdfstring{$\kappa$}{kappa} 等于生成函数的有理次数与 Hankel 秩]\label{cor:xi-scan-fingerprint-mcmillan-degree-hankel-rank}
\leanverified{paper\_xi\_scan\_fingerprint\_mcmillan\_degree\_hankel\_rank}
令生成函数
\[
U(z):=\sum_{n\ge 0}u_n z^n.
\]
则由 \eqref{eq:xi-scan-fingerprint-one-sided-exponential-sum} 得到有理表示
\[
\boxed{\;
U(z)=\sum_{j=1}^{\kappa}\frac{c_j}{1-r_j z}\;}
\]
从而 \(U\) 的最小分母次数（McMillan 度）在最小表示下等于 \(\kappa\)（重根情形按重数计）。
进一步，对块 Hankel 矩阵
\(
\mathsf{H}_N:=(u_{p+q})_{0\le p,q\le N}
\)
存在 \(N_0\) 使对所有 \(N\ge N_0\) 有 \(\mathrm{rank}(\mathsf{H}_N)=\kappa\)（见结论 \ref{con:xi-scan-hankel-rank-equals-defect}）。
因此缺陷度 \(\kappa\) 同时等于：
(i) 最小有理实现度；
(ii) 最小收缩实现维数（命题 \ref{prop:xi-scan-fingerprint-minimal-contraction-realization}）；
(iii) Hankel 秩稳定值。
\end{corollary}
\begin{proof}[证明要点]
有理表示由几何级数展开 \(1/(1-rz)=\sum_{n\ge 0}r^n z^n\) 逐项求和得到。
分母次数与 \(\kappa\) 的一致性在最小表示下来自极点个数（按重数计）。
Hankel 秩结论已由 \(\sum_j c_j r_j^n\) 的最小 realization 理论给出（定理 \ref{thm:hankel-rank-minimal} 与结论 \ref{con:xi-scan-hankel-rank-equals-defect}）。证毕。
\end{proof}

\begin{proposition}[深度--高度的张量分解唯一性：模长轨道与相位载体的不可替代直积]\label{prop:xi-scan-depth-height-tensor-factorization}
\leanverified{paper\_xi\_scan\_depth\_height\_tensor\_factorization}
写 \(r_j=\rho_j e^{-\mathrm{i}\theta_j}\)（\(\rho_j:=|r_j|\in(0,1)\)）。
则谱指纹分解为
\[
u_n=\sum_{j=1}^{\kappa}c_j\,\rho_j^{n}\,e^{-\mathrm{i}n\theta_j}.
\]
定义纯深度模长轨道
\[
v_n:=\sum_{j=1}^{\kappa}c_j\,\rho_j^{n}\qquad(n\ge 0).
\]
则 \(\{(\rho_j,c_j)\}\)（等价于任意一变量深度轨道）只能恢复模长衰减结构，而相位载体 \(\{e^{-\mathrm{i}n\theta_j}\}\) 属于不可替代的独立输入：不存在仅凭深度侧轨道而内生恢复相位 \(\theta\) 的制度内通道；任何试图消除此直积分裂的构造必在可审计闭包下退化为同余折叠或不可辨识回收。
该结论是命题 \ref{prop:xi-depth-height-minimal-coupling} 在离散地址谱指纹口径下的直接特化。
\end{proposition}
\begin{proof}
由 \(|r_j|=\rho_j\) 可知 \((v_n)\) 等价编码 \(\{(\rho_j,c_j)\}\) 的幅度谱；而 \(\theta_j\) 仅以字符 \(n\mapsto e^{-\mathrm{i}n\theta_j}\) 进入。
命题 \ref{prop:xi-depth-height-minimal-coupling} 断言：幅度谱与相位谱在制度上形成最小直积耦合，高度（相位）必须通过独立载体外置；将其中 \(e^{-\mathrm{i}\gamma_j n\Delta}\) 的角色限制到格点字符即得本命题。证毕。
\end{proof}

\begin{proposition}[“系数 4”在谱指纹侧的二重 Cayley 起源：\texorpdfstring{$2$}{2} 的两次规范化合成]\label{prop:xi-scan-coefficient-four-double-cayley}
\leanverified{paper\_xi\_scan\_coefficient\_four\_double\_cayley}
在扫描谱指纹的 Carath\'eodory 提升中，基本核满足
\[
\frac{1+w}{1-w}=1+2\sum_{n\ge 1}w^n\qquad(|w|<1),
\]
故 \(\mathscr{C}_u(z)\) 的系数前因子严格为 \(2\)（\eqref{eq:xi-scan-fingerprint-carath-series}）。
另一方面，在端点 \(-1\) 的圆盘化/半平面化规一化中，Cayley 变换的共形因子为 \(2\)，其平方作为密度拉回常数在边界 Poisson/Busemann 坐标中严格出现为 \(4\)（结论 \ref{con:xi-four-from-cayley-conformal-factor}）。
因此在“谱指纹 \(\to\) Carath\'eodory \(\to\) 边界密度”复合链路中，所有出现的 \(2\)/\(4\) 常数均来自规范化封口，而不承载新的动力学自由度。
\end{proposition}
\begin{proof}
第一部分由 Cayley 变换的幂级数展开直接给出。
第二部分已由结论 \ref{con:xi-four-from-cayley-conformal-factor} 证明：\(4\) 是共形因子 \(2\) 的平方，而在可审计不变量中仅以统一重标定方式出现并可在比值坐标中消去。
合并二者即得结论。证毕。
\end{proof}

\paragraph{纯深度子类：低阶 Hankel 判别式链与逐级源数下界}
以下把相位载体回收为实轴模长：假设 \(r_j=\rho_j\in(0,1)\subset\RR\) 并保持 \(c_j>0\)，从而 \((u_n)\) 为 \((0,1)\) 上有限正测度的幂矩序列。

\begin{proposition}[双源门槛的最小代数证书：\texorpdfstring{$\kappa\ge 2$}{kappa>=2} 的首个判别式]\label{prop:xi-scan-two-source-min-discriminant}
\leanverified{paper\_xi\_scan\_two\_source\_min\_discriminant}
在纯深度子类中，定义二阶 Hankel 判别式
\[
\Delta_1
:=\det\begin{pmatrix}u_0 & u_1\\[2pt] u_1 & u_2\end{pmatrix}
=u_0u_2-u_1^2.
\]
则必有 \(\Delta_1\ge 0\)，并且存在显式分解
\[
\boxed{\;
\Delta_1=\sum_{1\le j<k\le \kappa}c_jc_k(\rho_j-\rho_k)^2\; }.
\]
从而在最小表示口径下（\(\rho_j\) 两两不同），下列命题等价：
\[
\boxed{\;\kappa\ge 2\quad\Longleftrightarrow\quad \Delta_1>0.\;}
\]
\end{proposition}
\begin{proof}
展开得
\[
u_0u_2-u_1^2
=\Bigl(\sum_j c_j\Bigr)\Bigl(\sum_k c_k\rho_k^2\Bigr)-\Bigl(\sum_j c_j\rho_j\Bigr)^2
=\sum_{j,k}c_jc_k(\rho_k^2-\rho_j\rho_k).
\]
对称化并整理为 \(\sum_{j<k}c_jc_k(\rho_j-\rho_k)^2\) 即得分解，因而 \(\Delta_1\ge 0\)，且在 \(\rho_j\) 两两不同且 \(\kappa\ge 2\) 时严格为正。
反之若 \(\Delta_1=0\)，则上式强制 \(\rho_j\) 全相等，从而最小表示退化为单源。证毕。
\end{proof}

\begin{proposition}[双源持续性与证书稳定性：二源判别的条件数由谱隙唯一控制]\label{prop:xi-scan-two-source-stability-gap}
\leanverified{paper\_xi\_scan\_two\_source\_stability\_gap}
仍在纯深度子类中，取排序使 \(\rho_1>\rho_2\ge\cdots\ge\rho_\kappa\)，并定义谱隙比
\[
\eta:=\frac{\rho_2}{\rho_1}\in(0,1).
\]
则在固定总质量 \(u_0=\sum_j c_j\) 且权重不退化（例如 \(c_1,c_2\) 有一致正下界）的参数域内，
任意将“二源”判别式 \(\Delta_1\) 从正值推向 \(0\) 的扰动都不可避免地经过由 \((1-\eta)\) 控制的病态区；
更定量地，对二源分离量 \(\rho_1-\rho_2=\rho_1(1-\eta)\) 的局部反演，其 Lipschitz 条件数至少按 \((1-\eta)^{-1}\) 量级发散。
\end{proposition}
\begin{proof}[证明要点]
由命题 \ref{prop:xi-scan-two-source-min-discriminant}，有
\(
\Delta_1\ge c_1c_2(\rho_1-\rho_2)^2
\)。
故在 \(c_1,c_2\) 具有一致正下界时，
\(
\rho_1-\rho_2\le C\sqrt{\Delta_1}
\)
且反向导数尺度满足
\(
\mathrm{d}(\rho_1-\rho_2)\sim \Delta_1^{-1/2}\mathrm{d}\Delta_1
\)；
当 \(\rho_2\uparrow\rho_1\)（即 \(1-\eta\downarrow 0\)）时必触发 \(\Delta_1\downarrow 0\) 并导致条件数至少按 \((\rho_1-\rho_2)^{-1}\asymp (1-\eta)^{-1}\) 发散。证毕。
\end{proof}

\begin{proposition}[从双源到多源：低阶判别式链给出 \texorpdfstring{$\kappa$}{kappa} 的逐级下界]\label{prop:xi-scan-hankel-determinant-chain-kappa}
\leanverified{paper\_xi\_scan\_hankel\_determinant\_chain\_kappa}
在纯深度子类中，对 \(m\ge 0\) 定义 Hankel 判别式链
\[
\Delta_m:=\det(u_{p+q})_{0\le p,q\le m}.
\]
则对任意 \(m\) 有 \(\Delta_m\ge 0\)。
并且在最小表示口径下（\(\rho_j\) 两两不同），成立逐级等价：
\[
\boxed{\;
\kappa\ge m+1\quad\Longleftrightarrow\quad \Delta_m>0,\qquad
\kappa\le m\quad\Longleftrightarrow\quad \Delta_m=0.\;}
\]
因此 \(\{\Delta_m\}\) 构成一条完全局部的“源数下界证书链”：无需反演 \(\{\rho_j\}\)，仅需逐级检验有限阶 Hankel 行列式的严格正性，即可在制度可见域内确定“至少 \(m+1\) 源”的事实；其失效（\(\Delta_m=0\)）即对应第 \(m+1\) 个自由度被结构性回收。
\end{proposition}
\begin{proof}[证明要点]
记向量 \(v(\rho):=(1,\rho,\dots,\rho^{m})^\top\in\RR^{m+1}\)。
则 Hankel 块可写为 Gram 型分解
\[
(u_{p+q})_{0\le p,q\le m}
=\sum_{j=1}^{\kappa}c_j\,v(\rho_j)v(\rho_j)^\top\succeq 0,
\]
从而 \(\Delta_m\ge 0\)。
若 \(\kappa\le m\)，则上述矩阵秩不超过 \(\kappa\)，故 \(\Delta_m=0\)。
若 \(\kappa\ge m+1\) 且 \(\rho_j\) 两两不同，则存在 \(m+1\) 个节点使对应的 Vandermonde 向量组线性无关，Gram 矩阵满秩，从而 \(\Delta_m>0\)。
等价亦可由显式 Vandermonde 展开给出：
\[
\Delta_m=\sum_{1\le j_0<\cdots<j_m\le \kappa}
\Bigl(\prod_{\ell=0}^{m}c_{j_\ell}\Bigr)
\prod_{0\le p<q\le m}(\rho_{j_q}-\rho_{j_p})^2,
\]
其每一项非负且在存在 \(m+1\) 个互异节点时至少一项严格为正。证毕。
\end{proof}

\endinput
